\[ %%%%%%%%%%%%%%%%%%%%%%%%%%%%%%%%%%%%%%%%%%%%%%%%%%%%%%%%%%%%%%%%%%%%%%%%%%%%%%%% \subsection{Why Physics Guidance Instead of Physics-Informed Learning?} %%%%%%%%%%%%%%%%%%%%%%%%%%%%%%%%%%%%%%%%%%%%%%%%%%%%%%%%%%%%%%%%%%%%%%%%%%%%%%%% The integration of physical knowledge into machine learning has become one of the defining directions of modern Scientific Machine Learning. Among the numerous approaches proposed in recent years, Physics-Informed Neural Networks (PINNs) have established an influential paradigm by incorporating residuals of governing differential equations directly into the optimization objective. This strategy has demonstrated that first-principles physical models can effectively regularize data-driven learning while reducing dependence on large annotated datasets. However, the concept of \emph{physics-informed learning} is primarily associated with a specific mechanism of integrating physics into model training. In most existing formulations, physical knowledge enters the optimization process through additional loss terms that penalize violations of governing equations, boundary conditions, or initial conditions. Consequently, physics acts mainly as an external source of supervision during parameter optimization, while the internal computational architecture of the neural network remains largely unchanged. Although this formulation has achieved remarkable success across numerous applications, it represents only one possible strategy for incorporating physical knowledge into machine learning. Physical reasoning can influence computational inference at multiple levels beyond loss-function regularization. For example, physical principles may determine graph topology, constrain information propagation, regulate latent state evolution, control numerical stability, guide adaptive message passing, preserve conservation laws, influence uncertainty estimation, and determine admissible state transitions throughout recursive prediction. These observations motivate the broader concept of \emph{Physics-Guided Learning}. Within this perspective, physical knowledge is not viewed solely as an auxiliary optimization constraint but rather as a fundamental computational principle governing the behavior of the learning architecture itself. Instead of asking whether predicted states satisfy governing equations only after they have been generated, a physics-guided architecture continuously incorporates physical reasoning throughout the entire computational pipeline, including representation learning, graph construction, message passing, temporal evolution, latent dynamics, stability control, and uncertainty management. From this viewpoint, Physics-Informed Learning may be regarded as a particular instance of the more general Physics-Guided paradigm, where guidance is implemented primarily through equation-based loss regularization. Physics-Guided Learning, in contrast, encompasses a broader family of computational mechanisms capable of embedding physical knowledge into multiple architectural components simultaneously. These mechanisms are complementary rather than mutually exclusive. Indeed, residual equation losses naturally remain an important element within a physics-guided framework, but they no longer constitute its sole source of physical reasoning. The PHYSGEN framework proposed in this work adopts this broader architectural perspective. Rather than introducing physics only through additional optimization objectives, PHYSGEN distributes physical guidance across several interacting computational modules. Physical principles influence graph construction, message-passing dynamics, latent state propagation, conservation-aware regularization, adaptive stability control, and recursive prediction. As a consequence, physical consistency becomes an intrinsic property of the computational architecture rather than a post hoc correction applied during optimization. This architectural interpretation provides greater flexibility for extending scientific machine learning to heterogeneous physical domains involving irregular computational meshes, evolving graph topologies, coupled multi-physics systems, and long-horizon autoregressive simulations. Furthermore, it naturally accommodates hybrid computational strategies in which analytical models, numerical solvers, experimental observations, and data-driven neural representations cooperate within a unified differentiable framework. Accordingly, throughout this paper the term \emph{Physics-Guided Graph Networks} refers to graph neural architectures whose computational behavior is continuously directed by physical principles across multiple levels of inference, optimization, and temporal evolution, thereby extending conventional equation-constrained learning toward a more comprehensive computational paradigm for scientific intelligence. \] 
